\clearpage
\section{测试计划与用例设计}
\subsection{测试策略与执行条件}
\label{sec:impl-13-plan}
\par\medskip\noindent\textbf{测试目标与优先顺序。}
本轮目标是检查当前实现的关键规则、账号隔离、错误恢复和异步界面，不以覆盖截图数量代替功能覆盖。按风险排序，先测认证与越权，再测会话/抽取一致性，随后测模型适配、媒体字节与页面竞争；报告排版检查作为交付检查单独执行。

\begin{table}[H]\centering\small
\caption{测试层次、方法与验证对象}
\begin{tabularx}{\linewidth}{@{}p{22mm}p{30mm}X@{}}\toprule
层次 & 方法 & 主要检查对象 \\ \midrule
规则单元 & 等价类、边界值 & 长度、分数、时间、标签与状态迁移 \\
数据库 & 约束、回滚、重开 & 唯一性、事务失败、迁移幂等与持久化 \\
HTTP集成 & 角色决策表、异常注入 & 会话参与者、错误码、上游格式和限流 \\
浏览器 & 场景回归、竞争控制 & 草稿、风格、消息新旧、预览保存与视口 \\ \bottomrule
\end{tabularx}\end{table}

\par\medskip\noindent\textbf{夹具与执行条件。}
单元及接口测试使用测试账号、内存或临时SQLite，模型通过\DesignCode{fetchImpl}注入。固定时钟用于跨日和周期检查，固定随机票用于抽取边界，延迟Promise用于控制请求先后。每例结束清理服务、连接和数据库，不让一次失败污染其他用例。

准入条件为Node.js版本满足工程要求、源码能被加载、测试资源可创建且不依赖实际密钥。准出条件为计划执行的回归全部完成、失败有明确记录、关键权限或覆盖错误未遗留，以及报告中的数字可回到原始输出核对。没有执行的路径列为待测，不自动写成“通过”。

每个测试可包含多个断言；141是测试框架汇总项数，不是断言数、需求数或覆盖率。现有工程未生成语句/分支覆盖率报告，因此不能由全部用例通过推出“覆盖率100\%”。
